\addcontentsline{toc}{subsubsection}{Solution 40. The Monotonicity Ratchet}

\begin{dlsummary}[title={\bfseries Solution 40. The Monotonicity Ratchet},breakable]

\begin{proof}
The only solution is \(f(x)=x\). It satisfies both inequalities with equality.

\paragraph{1. A first consequence.} In the first inequality, \(f(x)>0\), so \(f(x+f(y))>y\). As \(x\) ranges over all positive reals, \(t=x+f(y)\) ranges over all \(t>f(y)\). Therefore
\[
t>f(y)\ \Longrightarrow\ f(t)>y. \tag{1}
\]

\paragraph{2. \(f\) is strictly increasing.} Take \(0<w<t\). By the second condition, \(f(f(w))\le w<t\). Apply (1) with \(y=f(w)\). Since \(t>f(f(w))\), we get
\[
f(t)>f(w).
\]

\paragraph{3. \(f(y)\le y\).} Suppose \(f(y)>y\) for some \(y\). Since \(f\) is strictly increasing, \(f(f(y))>f(y)>y\). This contradicts \(f(f(y))\le y\).

\paragraph{4. Collapse to a constant defect.} Using step 3 and the first condition,
\[
x+f(y)\ge f\bigl(x+f(y)\bigr)\ge f(x)+y.
\]
Rearranging gives
\[
f(y)-y\ge f(x)-x\qquad\text{for all }x,y>0.
\]
Swapping the roles of \(x\) and \(y\) gives the reverse inequality. Hence \(f(x)-x=-k\) is constant, and step 3 gives \(k\ge0\).

\paragraph{5. Fix the constant.} If \(k>0\), then \(f(k)=0\) with \(k>0\), which contradicts positivity of \(f\). So \(k=0\) and \(f(x)=x\). \(\blacksquare\)
\end{proof}
\end{dlsummary}

\begin{takeaways}
\begin{itemize}
\item \textbf{Derive monotonicity:} Substituting \(y=f(w)\) turned a composition bound into a statement about order.
\item \textbf{Use order against iteration:} For an increasing function, \(f(y)>y\) would force \(f(f(y))>y\), which contradicts an upper bound on \(f(f(y))\).
\item \textbf{Use the symmetry trick again:} Once \(f(y)-y\ge f(x)-x\) holds for all \(x,y\), swapping them gives equality.
\item \textbf{Let the domain finish the job:} Positivity of \(f\) removes the family \(f(x)=x-k\).
\end{itemize}
\end{takeaways}
